\documentclass[11pt]{article}
\usepackage[utf8]{inputenc}
\usepackage[T1]{fontenc}
\usepackage{lmodern}
\usepackage[spanish]{babel}
\usepackage{amsmath,amssymb}
\usepackage{graphicx}
\usepackage[left=2cm,right=2cm,top=2cm,bottom=2cm]{geometry}
\usepackage{enumitem}
\usepackage{booktabs}
\usepackage{array}
\usepackage{xcolor}
\usepackage{framed}

\spanishdecimal{.}

\setlist[enumerate,1]{label=\textbf{\arabic*.}, leftmargin=*, itemsep=6pt}
\setlist[enumerate,2]{label=(\alph*), itemsep=4pt}
\setlength{\parindent}{0pt}
\setlength{\parskip}{4pt}
\pagestyle{plain}

\newcommand{\Ein}{E_{\text{in}}}
\newcommand{\Eout}{E_{\text{out}}}
\newcommand{\Hyp}{\mathcal{H}}
\newcommand{\dvc}{d_{\text{VC}}}
\newcommand{\vx}{\mathbf{x}}
\newcommand{\vw}{\mathbf{w}}
\newcommand{\vy}{\mathbf{y}}
\newcommand{\R}{\mathbb{R}}
\newcommand{\T}{^{\top}}
\newcommand{\E}{\mathbb{E}}
\newcommand{\Prob}{\mathbb{P}}
\newcommand{\pts}[1]{\hfill\textbf{(#1~puntos)}}
\newcommand{\linea}[1]{\rule{#1}{0.4pt}}
\newcommand{\ptsinc}[1]{{\unskip\nobreak\hfil\penalty50\hskip1em\hbox{}\nobreak\hfill(#1)\parfillskip=0pt\finalhyphendemerits=0\par}}

\newcommand{\encabezado}[2]{%
\begin{center}
{\large\bfseries Reconocimiento de Patrones}\\[2pt]
Universidad de Sonora \textperiodcentered{} Licenciatura en Ciencias de la Computación\\
Semestre 2027-1 \textperiodcentered{} Prof.\ Julio Waissman Vilanova\\[10pt]
{\Large\bfseries #1}%
\ifdefined\tipodoc\\[4pt]{\large\itshape\tipodoc}\fi
\end{center}
\ifdefined\tipodoc\else
\vspace{6pt}
\textbf{Nombre:} \hrulefill\quad \textbf{Expediente:} \rule{3.5cm}{0.4pt}\par
\vspace{6pt}
\fi
#2
\par\medskip\hrule\medskip
}

\newcommand{\espacio}[1]{\par\vspace{#1}}
\newcommand{\nuevapagina}{\newpage}

\usepackage{tabularx}

\begin{document}

\encabezado{Examen intermedio para llevar \textperiodcentered{} Unidades 2 a 4}{%
\textbf{Plazo:} una semana a partir de su publicación. \textbf{Valor:} 100 puntos, 50 por problema.\par
\textbf{Reglas.} El trabajo es individual. Puedes consultar libros, apuntes, las libretas del curso y la documentación de las bibliotecas. Si usas asistentes de IA, indica en la entrega para qué los usaste; eres responsable de entender y poder explicar todo lo que entregas, y el profesor puede pedirte que expliques cualquier parte en una entrevista breve.\par
\textbf{Qué entregar.} (1) Un PDF con las derivaciones de los dos problemas, escrito en \LaTeX{} o en Markdown con fórmulas. (2) Una libreta por problema (\texttt{problema1.ipynb}, \texttt{problema2.ipynb}), \textbf{ejecutada de principio a fin} sin errores, que cargue los datos con \texttt{rdp.datos.cargar} y fije todas las semillas. Las gráficas deben tener ejes con nombre y unidades, y cada resultado debe ir acompañado de un comentario que diga qué significa.}

\begin{enumerate}

%% ===============================================================
\item \textbf{Pocas etiquetas, muchos pingüinos: EM semisupervisado.} \pts{50}

Medir el pico, la aleta y la masa de un pingüino es rápido; asegurar su especie (o su sexo) requiere a un especialista o una prueba genética. Supón que tenemos las medidas de todos los pingüinos, pero la etiqueta de solo unos cuantos. ¿Sirven de algo los pingüinos sin etiqueta?

La idea es usar el modelo generativo de la Unidad 2 con la información de la Unidad 4: una mezcla de $K$ gaussianas en la que \emph{cada componente es una clase}, $p(\vx) = \sum_k \pi_k\, \mathcal N(\vx \mid \boldsymbol\mu_k, \Sigma_k)$. Sea $L$ el conjunto de índices de los datos etiquetados $(\vx_i, y_i)$ y $U$ el de los datos sin etiqueta $\vx_i$, con $n = |L| + |U|$.

\begin{enumerate}
  \item \textbf{Derivación} (20 puntos).
  \begin{enumerate}[label=\arabic*.]
    \item Supón primero que los datos etiquetados se eligieron al azar entre todos. Escribe la log-verosimilitud $\ell(\theta)$ de los datos observados, con $\theta = \{\pi_k, \boldsymbol\mu_k, \Sigma_k\}$. \ptsinc{2}
    \item Deriva el algoritmo EM para maximizarla: di qué hace el paso E con los datos de $U$ y con los de $L$, y deriva el paso M para $\pi_k$, $\boldsymbol\mu_k$ y $\Sigma_k$. \ptsinc{6}
    \item Demuestra que $\ell(\theta)$ no disminuye de una iteración a la siguiente. Adapta el argumento de la cota inferior visto en clase y explica qué papel juegan los términos etiquetados. \ptsinc{4}
    \item Casos límite. (i) Si $U = \varnothing$, muestra que EM se detiene en un paso y que da los estimadores de QDA; con una covarianza compartida, $\Sigma_k = \Sigma$, los de LDA. (ii) Si $L = \varnothing$, ¿qué se pierde? \ptsinc{4}
    \item Con $m$ etiquetas por clase y $d = 4$ atributos, la covarianza de una clase estimada solo con sus etiquetados es singular si $m \le d$. Muestra que la previa $p(\Sigma_k) \propto \exp\!\big(-\tfrac{\varepsilon}{2}\operatorname{tr}\Sigma_k^{-1}\big)$, con $\varepsilon > 0$, cambia el paso M a $\Sigma_k = S_k + \tfrac{\varepsilon}{N_k} I$, donde $S_k$ es la covarianza ponderada de máxima verosimilitud, y explica por qué EM sigue siendo monótono, ahora para la log-posterior. \ptsinc{4}
  \end{enumerate}

  \item \textbf{Implementación} (15 puntos). Programa EM semisupervisado con NumPy (scikit-learn solo para partir, escalar y comparar), con covarianzas completas o compartidas y el término $\varepsilon$. Calcula el paso E en escala logarítmica, con \texttt{logsumexp}. Incluye pruebas automáticas (\texttt{assert}) de que: (i) con $U = \varnothing$ y $\varepsilon = 0$ se obtienen las medias y covarianzas de cada clase (QDA) y la covarianza combinada de \texttt{LinearDiscriminantAnalysis}; (ii) la log-posterior no baja en ninguna iteración.

  \item \textbf{Experimentos} (10 puntos). Usa \texttt{pinguinos} con las cuatro medidas numéricas (quita los renglones con faltantes). Parte una sola vez en 70\,\% para entrenamiento y 30\,\% para prueba, estratificando, y estandariza con el entrenamiento. Para $m \in \{2, 3, 5, 10, 20\}$ etiquetas por especie, repite 30 veces: elige al azar los $m$ etiquetados de cada especie; el resto del entrenamiento queda sin etiqueta. Compara la exactitud en prueba de: QDA y LDA entrenados solo con los etiquetados (con el mismo $\varepsilon$); EM semisupervisado con covarianzas completas y con covarianza compartida. Como referencias, QDA con todas las etiquetas y una mezcla sin etiquetas en la que cada componente toma la clase mayoritaria. Grafica la media y la desviación estándar contra $m$. Después repite todo con el \textbf{sexo} como clase ($K = 2$) y $m \in \{2, 5, 10, 20, 50\}$.

  \item \textbf{Análisis} (5 puntos). En no más de una página: ¿cuándo ayudaron los datos sin etiqueta y cuándo estorbaron? Relaciónalo con el supuesto del modelo generativo. Además, en el experimento las etiquetas se eligen $m$ por clase, no al azar: ¿qué parte de tu derivación deja de valer y cómo la corregiste?
\end{enumerate}

\nuevapagina
%% ===============================================================
\item \textbf{¿Cuántas bicicletas? Poisson regularizada y validación que respeta el tiempo.} \pts{50}

El sistema de bicicletas compartidas de Washington quiere pronosticar las rentas de cada hora para los próximos meses y planear la redistribución de bicicletas. Usarás \texttt{bicis\_washington} (2011--2012, una fila por hora). Para predecir \texttt{total} solo puedes usar información que se conocería antes de esa hora.

Modelo: $y \mid \vx \sim \operatorname{Poisson}(\mu)$ con $\ln\mu = \vw\T\vx$, donde $x_0 = 1$, y la penalización
\[
E(\vw) = \frac1N\sum_{n=1}^N \Big(e^{\vw\T\vx_n} - y_n\,\vw\T\vx_n\Big) + \frac{\lambda}{N}\sum_{j \ge 1} w_j^2 ,
\]
que no penaliza el intercepto $w_0$.

\begin{enumerate}
  \item \textbf{Derivación} (15 puntos).
  \begin{enumerate}[label=\arabic*.]
    \item Muestra que el primer término de $E(\vw)$ es, salvo una constante y el factor $1/N$, el negativo de la log-verosimilitud. Calcula el gradiente $\nabla E$ y la hessiana $\nabla^2 E$ en forma matricial. \ptsinc{4}
    \item Demuestra que, con $\lambda > 0$, $E$ es estrictamente convexa (aunque $w_0$ no se penalice), y concluye que su mínimo, si existe, es único. \ptsinc{4}
    \item Muestra que el paso de Newton se puede escribir como un problema de mínimos cuadrados ponderados con una respuesta de trabajo $\mathbf z$ (IRLS). Identifica los pesos y $\mathbf z$. \ptsinc{4}
    \item La clase \texttt{PoissonRegressor} de scikit-learn minimiza $\frac{1}{2N}\sum_n D(y_n, \mu_n) + \frac{\alpha}{2}\sum_{j \ge 1} w_j^2$, con la devianza $D(y, \mu) = 2\big(y\ln\frac{y}{\mu} - y + \mu\big)$. ¿Qué $\alpha$ da el mismo mínimo que $E$? \ptsinc{3}
  \end{enumerate}

  \item \textbf{Implementación} (15 puntos). Programa el método de Newton con búsqueda lineal por retroceso (empieza en $w_0 = \ln\bar y$ y los demás en cero) y verifica con un \texttt{assert} que tus coeficientes coinciden con los de \texttt{PoissonRegressor} con el $\alpha$ del inciso anterior. Construye los atributos con, al menos: \emph{one-hot} de la hora combinada con \texttt{laboral} (el patrón de un día laboral no es el de un domingo), \emph{one-hot} de \texttt{clima} y de \texttt{mes}, un indicador del año 2012, \texttt{temp}, $\texttt{temp}^2$, \texttt{humedad} y \texttt{viento}, estandarizados con el entrenamiento. Puedes agregar otros si los justificas. Revisa con cuidado qué columnas \emph{no} puedes usar.

  \item \textbf{Validación} (15 puntos). Aparta como prueba del 1 de octubre al 31 de diciembre de 2012 y no la uses hasta el final. Con el resto, elige $\lambda$ en una rejilla logarítmica (por ejemplo, de $10^{-2}$ a $10^{5}$) con dos esquemas: (i) validación cruzada de 5 particiones con las horas revueltas al azar; (ii) validación hacia adelante por mes: para cada mes de enero a septiembre de 2012, entrena con todos los meses anteriores y valida con ese mes. Usa como error la devianza media de Poisson (\texttt{mean\_poisson\_deviance}). Presenta una tabla con el error estimado por cada esquema para cada $\lambda$, el $\lambda$ que elige cada uno y la devianza en prueba de los dos modelos elegidos. Compara además, en prueba, con dos referencias: la media de rentas por hora y tipo de día (\texttt{laboral}) en el entrenamiento, y una regresión lineal de $\ln(1 + y)$ con los mismos atributos, regresada a la escala original.

  \item \textbf{Análisis} (5 puntos). En no más de una página: ¿por qué difieren tanto las estimaciones de los dos esquemas? ¿Cuál es honesta para el uso que se le dará al modelo? ¿Qué modelo recomendarías y con qué error esperado? Menciona al menos una limitación de tu evaluación.
\end{enumerate}

\end{enumerate}

\nuevapagina
\section*{Criterios de evaluación}

\begin{center}
\renewcommand{\arraystretch}{1.25}
\begin{tabularx}{\textwidth}{>{\raggedright\arraybackslash}p{4.2cm} c X}
\toprule
Componente & Puntos & Qué se espera \\
\midrule
\multicolumn{3}{l}{\textbf{Problema 1. EM semisupervisado}} \\
Derivación & 20 & Verosimilitud correcta, pasos E y M derivados (no solo enunciados), prueba de monotonía, casos límite y regularización. \\
Implementación & 15 & EM con NumPy, numéricamente estable, con las dos pruebas automáticas pasando. \\
Experimentos & 10 & Protocolo descrito y reproducible, las comparaciones pedidas para especie y sexo, gráficas legibles. \\
Análisis & 5 & Conclusiones apoyadas en los resultados propios; discusión del muestreo por clase. \\
\midrule
\multicolumn{3}{l}{\textbf{Problema 2. Poisson y validación temporal}} \\
Derivación & 15 & Gradiente y hessiana, convexidad estricta con el intercepto sin penalizar, IRLS y relación con scikit-learn. \\
Implementación & 15 & Newton con retroceso verificado contra scikit-learn; atributos legítimos y justificados. \\
Validación & 15 & Prueba apartada, dos esquemas de validación bien hechos, tabla y referencias. \\
Análisis & 5 & Por qué difieren los esquemas, cuál es honesto, recomendación y limitaciones. \\
\bottomrule
\end{tabularx}
\end{center}

\textbf{Descuentos generales.} Libreta que no corre de principio a fin: $-10$ por problema. Resultados sin comentario: hasta $-5$ por problema. Gráficas sin ejes con nombre: $-1$ cada una. Derivaciones copiadas sin entender (detectado en la entrevista): el inciso vale cero.


\end{document}
